

\begin{problem}[Matrix Systems and Polynomial Curve Fitting]
A quadratic function has a rule of the form $f(x) = ax^2 + bx + c$. It is known that the points $(-1, 12)$ and $(1, 4)$ are on the graph of $f$, alongside a third point $(2, 9)$.

\begin{enumerate}[label=\textbf{(\alph*)}]
    \item By substituting the known points, write down three linear equations satisfied by $a, b$, and $c$. Formulate this information as a matrix equation of the form $M \mathbf{v} = \mathbf{k}$, where $\mathbf{v} = \begin{bmatrix} a \\ b \\ c \end{bmatrix}$.
    \item The inverse of the coefficient matrix $M$ is given by $M^{-1} = \frac{1}{6} \begin{bmatrix} 1 & -3 & 2 \\ -3 & 3 & 0 \\ 2 & 6 & -2 \end{bmatrix}$. Use this inverse to solve the matrix equation and determine the exact values of $a, b$, and $c$.
    \item Suppose the $y$-coordinate of the third point is changed to a parameter $p$, such that the point is now $(2, p)$. Write down the new constant vector $\mathbf{k}$, and use $M^{-1}$ to determine the coefficients $a, b$, and $c$ in terms of $p$.
    \item Determine the specific value of $p$ for which the quadratic function degenerates into a linear function (i.e., $a = 0$). Geometrically, what does this imply about the arrangement of the three points on the Cartesian plane?
\end{enumerate}
\end{problem}

\begin{hint}
\begin{itemize}
    \item \textbf{Part (a):} Substitute the $x$ and $y$ values of each coordinate pair into the quadratic equation to generate three distinct linear equations. The coefficients of $a, b,$ and $c$ will form your $3 \times 3$ matrix $M$.
    \item \textbf{Part (b):} To solve $M \mathbf{v} = \mathbf{k}$, pre-multiply both sides by the inverse matrix so that $\mathbf{v} = M^{-1} \mathbf{k}$.
    \item \textbf{Part (c):} Update the third entry of your $\mathbf{k}$ vector to $p$. Multiply the matrix $M^{-1}$ by this new vector using standard row-by-column multiplication.
    \item \textbf{Part (d):} Take your parametric expression for $a$ from part (c) and set it equal to 0. If a quadratic curve becomes a straight line, consider what that means for the points that define it.
\end{itemize}
\end{hint}

\begin{solution}
\textbf{(a)} \\
Substituting $(-1, 12)$: $a(-1)^2 + b(-1) + c = 12 \implies a - b + c = 12$ \\
Substituting $(1, 4)$: $a(1)^2 + b(1) + c = 4 \implies a + b + c = 4$ \\
Substituting $(2, 9)$: $a(2)^2 + b(2) + c = 9 \implies 4a + 2b + c = 9$ \\
Matrix equation: $\begin{bmatrix} 1 & -1 & 1 \\ 1 & 1 & 1 \\ 4 & 2 & 1 \end{bmatrix} \begin{bmatrix} a \\ b \\ c \end{bmatrix} = \begin{bmatrix} 12 \\ 4 \\ 9 \end{bmatrix}$.

\medskip
\textbf{(b)} \\
$\mathbf{v} = M^{-1} \mathbf{k} = \frac{1}{6} \begin{bmatrix} 1 & -3 & 2 \\ -3 & 3 & 0 \\ 2 & 6 & -2 \end{bmatrix} \begin{bmatrix} 12 \\ 4 \\ 9 \end{bmatrix}$ \\
$a = \frac{1}{6} (1(12) - 3(4) + 2(9)) = \frac{1}{6}(12 - 12 + 18) = 3$ \\
$b = \frac{1}{6} (-3(12) + 3(4) + 0(9)) = \frac{1}{6}(-36 + 12) = -4$ \\
$c = \frac{1}{6} (2(12) + 6(4) - 2(9)) = \frac{1}{6}(24 + 24 - 18) = 5$ \\
Therefore, $a = 3$, $b = -4$, and $c = 5$.

\medskip
\textbf{(c)} \\
With the new point $(2, p)$, the constant vector becomes $\mathbf{k} = \begin{bmatrix} 12 \\ 4 \\ p \end{bmatrix}$. \\
$\mathbf{v} = \frac{1}{6} \begin{bmatrix} 1 & -3 & 2 \\ -3 & 3 & 0 \\ 2 & 6 & -2 \end{bmatrix} \begin{bmatrix} 12 \\ 4 \\ p \end{bmatrix}$ \\
$a = \frac{1}{6} (12 - 12 + 2p) = \frac{2p}{6} = \frac{p}{3}$ \\
$b = \frac{1}{6} (-36 + 12 + 0) = -4$ \\
$c = \frac{1}{6} (24 + 24 - 2p) = \frac{48 - 2p}{6} = 8 - \frac{p}{3}$.

\medskip
\textbf{(d)} \\
For the function to be linear, $a = 0$. \\
$\frac{p}{3} = 0 \implies p = 0$. \\
If $a = 0$, the curve is a straight line $y = -4x + 8$. Geometrically, this implies that the three points $(-1, 12)$, $(1, 4)$, and $(2, 0)$ are perfectly collinear.
\end{solution}

\begin{takeaways}
\item \textbf{Vandermonde Matrices:} The coefficient matrix formed by fitting a polynomial to coordinate points is known as a Vandermonde matrix. Its determinant is always non-zero as long as the $x$-coordinates are distinct, guaranteeing a unique polynomial fit.
\item \textbf{Parametric Matrix Operations:} Treating unknowns as parameters (like $p$) within a matrix equation allows you to generalize solutions for an entire family of curves at once, rather than re-solving the system from scratch.
\item \textbf{Geometric Limits:} Setting the leading coefficient of a polynomial to zero algebraically represents flattening the curve's curvature, visually forcing the defining points into a straight line.
\end{takeaways}
